% Bagian: proof-theory
% Bab: sequent-calculus
% Subbagian: translations

\documentclass[../../../include/open-logic-section]{subfiles}

\begin{document}

\olfileid[id]{pt}{seq}{tr}

\olsection{Penerjemahan antara \Log{G1c} dan \Log{G3c}}

Kita telah menyebutkan bahwa \Log{G1c} dan \Log{G3c} sama-sama sahih dan
lengkap bagi logika klasik, yakni masing-masing membuktikan tepat sekuen-sekuen
yang benar dalam setiap !!{structure}. Secara khusus, keduanya membuktikan
sekuen-sekuen yang sama. Sekarang kita dapat membuktikan fakta ini tanpa
memutar melalui teorema kesahihan dan kelengkapan, melainkan dengan cara yang
murni bersifat teori bukti. Bukti semacam itu memberikan
\emph{transformasi} dari !!{proof}s dalam \Log{G1c} menjadi !!{proof}s dalam
\Log{G3c}, dan sebaliknya.

\begin{prop}
Jika $\Log{G1c} \Proves S$, maka $\Log{G3c} \Proves S$.
\end{prop}

\begin{proof}
  Kita tunjukkan bahwa jika $\pi$ merupakan !!a{proof} dalam \Log{G1c}, maka
  terdapat !!a{proof}~$\pi'$ bagi~$S$ dalam \Log{G3c}. Kita melakukannya
  dengan induksi pada $n = \pheight{\pi}$.

  Jika $n = 0$, maka $\pi$ merupakan aksioma. Aksioma $\lfalse \Sequent
  \quad$ juga merupakan aksioma dalam~\Log{G3c} (ambil konteks $\Gamma,
  \Delta$ kosong). Dalam \Log{G1c}, aksioma $!A \Sequent !A$ diperbolehkan
  untuk $!A$ sembarang, bukan hanya yang atomik. Namun, dalam
  \olref[ex]{prop:G3c-axioms} kita telah membuktikan bahwa semua sekuen
  semacam itu !!{provable} dalam~\Log{G3c}.
  
  Jika $n>0$, $\pi$ berakhir pada suatu aturan~$R$. Hipotesis induksi menjamin
  bahwa terdapat !!{proof}s bagi premis-premis aturan itu dalam~\Log{G3c},
  yakni sub-!!{proof}s langsungnya memiliki terjemahan dalam~\Log{G3c}. Jika
  aturan~$R$ juga merupakan aturan dalam \Log{G3c}, kita menerapkannya pada
  !!{proof}s terjemahan bagi premis-premis tersebut. Untuk aturan-aturan yang
  bukan merupakan aturan dalam~\Log{G3c}, kita telah menunjukkan bahwa
  aturan-aturan itu setidaknya admisibel. Jadi, jika aturan~$R$ adalah
  $\LeftR{\land}$ atau $\RightR{\lor}$, terapkan
  \olref[adm]{prop:lor-G3-adm}. Jika aturan itu aturan pelemahan, terapkan
  \olref[adm]{prop:weak-G3c-adm}. Jika aturan itu aturan kontraksi, terapkan
  \olref[inv]{prop:G3c-cont-adm}.
\end{proof}

\begin{prop}
Jika $\Log{G3c} \Proves S$, maka $\Log{G1c} \Proves S$.
\end{prop}

\begin{proof}
  Sekali lagi kita menggunakan induksi pada !!{height} suatu !!a{proof}~$\pi$
  bagi $S$. Dalam~\Log{G1c}, kita dapat !!{prove} setiap aksioma \Log{G3c}
  dari suatu aksioma beserta beberapa penerapan aturan \LeftR{\Weakening} dan
  \RightR{\Weakening}:
  \[
  \Axiom$!C \fCenter !C$
  \RightLabel{\LeftR{\Weakening}}
  \Deduce$!C, \Gamma \fCenter !C$
  \RightLabel{\RightR{\Weakening}}
  \Deduce$!C, \Gamma \fCenter \Delta, !C$
  \DisplayProof
\qquad
  \Axiom$\lfalse \fCenter$
  \RightLabel{\LeftR{\Weakening}}
  \Deduce$\lfalse, \Gamma \fCenter$
  \RightLabel{\RightR{\Weakening}}
  \Deduce$\lfalse, \Gamma \fCenter \Delta$
  \DisplayProof
  \]
  Jika $\pi$ berakhir pada aturan~$R$, hipotesis induksi menghasilkan
  !!{proof}s dalam \Log{G1c} bagi premis-premisnya, dan kita dapat menerapkan
  aturan~$R$ yang sama pada !!{proof}s tersebut. Pengecualiannya ialah aturan
  \LeftR{\land} dan \RightR{\lor}, yang berbeda antara \Log{G1c}
  dan~\Log{G3c}. Versi \Log{G3c} dapat diturunkan dalam \Log{G1c} sebagai
  berikut:
  \[
  \Axiom$!A, !B, \Gamma \fCenter \Delta$
  \RightLabel{\LeftR{\land}}
  \UnaryInf$!A \land !B, !B, \Gamma \fCenter \Delta$
  \RightLabel{\LeftR{\land}}
  \UnaryInf$!A \land !B, !A \land !B, \Gamma \fCenter \Delta$
  \RightLabel{\LeftR{\Contraction}}
  \UnaryInf$!A \land !B, \Gamma \fCenter \Delta$
    \DisplayProof
\qquad
  \Axiom$\Gamma \fCenter \Delta, !A, !B$
  \RightLabel{\RightR{\lor}}
  \UnaryInf$\Gamma \fCenter \Delta, !A \lor !B, !B$
  \RightLabel{\RightR{\lor}}
  \UnaryInf$\Gamma \fCenter \Delta, !A \lor !B, !A \lor !B$
  \RightLabel{\RightR{\Contraction}}
  \UnaryInf$\Gamma \fCenter \Delta, !A \lor !B$
    \DisplayProof
  \]
\end{proof}

\end{document}
